% Options for packages loaded elsewhere
\PassOptionsToPackage{unicode}{hyperref}
\PassOptionsToPackage{hyphens}{url}
\documentclass[
  ignorenonframetext,
]{beamer}
\newif\ifbibliography
\usepackage{pgfpages}
\setbeamertemplate{caption}[numbered]
\setbeamertemplate{caption label separator}{: }
\setbeamercolor{caption name}{fg=normal text.fg}
\beamertemplatenavigationsymbolsempty
% remove section numbering
\setbeamertemplate{part page}{
  \centering
  \begin{beamercolorbox}[sep=16pt,center]{part title}
    \usebeamerfont{part title}\insertpart\par
  \end{beamercolorbox}
}
\setbeamertemplate{section page}{
  \centering
  \begin{beamercolorbox}[sep=12pt,center]{section title}
    \usebeamerfont{section title}\insertsection\par
  \end{beamercolorbox}
}
\setbeamertemplate{subsection page}{
  \centering
  \begin{beamercolorbox}[sep=8pt,center]{subsection title}
    \usebeamerfont{subsection title}\insertsubsection\par
  \end{beamercolorbox}
}
% Prevent slide breaks in the middle of a paragraph
\widowpenalties 1 10000
\raggedbottom
\AtBeginPart{
  \frame{\partpage}
}
\AtBeginSection{
  \ifbibliography
  \else
    \frame{\sectionpage}
  \fi
}
\AtBeginSubsection{
  \frame{\subsectionpage}
}
\usepackage{iftex}
\ifPDFTeX
  \usepackage[T1]{fontenc}
  \usepackage[utf8]{inputenc}
  \usepackage{textcomp} % provide euro and other symbols
\else % if luatex or xetex
  \usepackage{unicode-math} % this also loads fontspec
  \defaultfontfeatures{Scale=MatchLowercase}
  \defaultfontfeatures[\rmfamily]{Ligatures=TeX,Scale=1}
\fi
\usepackage{lmodern}
\ifPDFTeX\else
  % xetex/luatex font selection
\fi
% Use upquote if available, for straight quotes in verbatim environments
\IfFileExists{upquote.sty}{\usepackage{upquote}}{}
\IfFileExists{microtype.sty}{% use microtype if available
  \usepackage[]{microtype}
  \UseMicrotypeSet[protrusion]{basicmath} % disable protrusion for tt fonts
}{}
\makeatletter
\@ifundefined{KOMAClassName}{% if non-KOMA class
  \IfFileExists{parskip.sty}{%
    \usepackage{parskip}
  }{% else
    \setlength{\parindent}{0pt}
    \setlength{\parskip}{6pt plus 2pt minus 1pt}}
}{% if KOMA class
  \KOMAoptions{parskip=half}}
\makeatother
\usepackage{longtable,booktabs,array}
\newcounter{none} % for unnumbered tables
\usepackage{calc} % for calculating minipage widths
\usepackage{caption}
% Make caption package work with longtable
\makeatletter
\def\fnum@table{\tablename~\thetable}
\makeatother
\setlength{\emergencystretch}{3em} % prevent overfull lines
\providecommand{\tightlist}{%
  \setlength{\itemsep}{0pt}\setlength{\parskip}{0pt}}
% Auto-generated by infrastructure/rendering/_pdf_latex_helpers.py::extract_math_font_preamble.
% Minimal header passed to Pandoc via -H so Beamer slides pick up the same math font
% as the combined-PDF path. Adjust the manuscript's preamble.md to change the font globally.
\usepackage{unicode-math}
\setmathfont{latinmodern-math.otf}
\usepackage{bookmark}
\IfFileExists{xurl.sty}{\usepackage{xurl}}{} % add URL line breaks if available
\urlstyle{same}
\hypersetup{
  hidelinks,
  pdfcreator={LaTeX via pandoc}}

\author{\texorpdfstring{}{}}
\date{}

\begin{document}

\section{Research Questions and Manuscript
Navigation}\label{sec:research-questions}

\begin{frame}[allowframebreaks]{Research Questions and Manuscript
Navigation}
This paper addresses eleven core research questions, each developed in
specific sections and yielding a principal contribution
(\textbf{C1}--\textbf{C11}, as enumerated in \autoref{sec:conclusion}).
The table below provides a navigational overview; the reading guide that
follows explains the logical dependencies between questions.

\begingroup
\small

\begin{longtable}[]{@{}
  >{\centering\arraybackslash}p{(\linewidth - 6\tabcolsep) * \real{0.0656}}
  >{\raggedright\arraybackslash}p{(\linewidth - 6\tabcolsep) * \real{0.7705}}
  >{\raggedright\arraybackslash}p{(\linewidth - 6\tabcolsep) * \real{0.0984}}
  >{\centering\arraybackslash}p{(\linewidth - 6\tabcolsep) * \real{0.0656}}@{}}
\caption{Core research questions mapped to corresponding manuscript
sections and resulting
contributions.}\label{tbl:research-questions}\tabularnewline
\toprule\noalign{}
\begin{minipage}[b]{\linewidth}\centering
\#
\end{minipage} & \begin{minipage}[b]{\linewidth}\raggedright
Research Question
\end{minipage} & \begin{minipage}[b]{\linewidth}\raggedright
Addressing Sections
\end{minipage} & \begin{minipage}[b]{\linewidth}\centering
C
\end{minipage} \\
\midrule\noalign{}
\endfirsthead
\toprule\noalign{}
\begin{minipage}[b]{\linewidth}\centering
\#
\end{minipage} & \begin{minipage}[b]{\linewidth}\raggedright
Research Question
\end{minipage} & \begin{minipage}[b]{\linewidth}\raggedright
Addressing Sections
\end{minipage} & \begin{minipage}[b]{\linewidth}\centering
C
\end{minipage} \\
\midrule\noalign{}
\endhead
\textbf{RQ1} & Can linguistic case systems---across all major
typological alignment patterns---be formalized within a single algebraic
framework, with roles as objects and grammatical relations as morphisms?
& \autoref{sec:case-systems}, \autoref{sec:case-categories} &
\textbf{C1} \\
\textbf{RQ2} & Do string diagrams derived from pregroup type reductions
provide a computationally and cognitively superior representation of
case derivations compared to linear notation? &
\autoref{sec:categorial-grammar}, \autoref{sec:case-type-logic},
\autoref{sec:syntactic-diagrams} & \textbf{C2} \\
\textbf{RQ3} & Can the DisCoCat meaning functor be extended with
case-typed noun spaces to unify formal and distributional
semantics---and does this unification generalize to discourse via
DisCoCirc? & \autoref{sec:categorical-semantics},
\autoref{sec:discocat-meaning-functor},
\autoref{sec:compact-closure-complexity},
\autoref{sec:discocirc-discourse} & \textbf{C3} \\
\textbf{RQ4} & Does equipping case categories with \([0,1]\)-enriched
hom-values yield a principled bridge between symbolic grammar and
statistical semantics, and can categorical magnitude serve as a
complexity invariant? & \autoref{sec:enriched-categories},
\autoref{sec:magnitude-homology} & \textbf{C4} \\
\textbf{RQ5} & Can classifying toposes and Morita equivalence (or
explicit bridge toposes) scaffold transfer of \textbf{topos-expressible
invariants} between distinct formalizations of case theory (typological,
type-logical, distributional, enriched), when equivalence is exhibited?
& \autoref{sec:topos-theory} & \textbf{C5} \\
\textbf{RQ6} & Are commutative diagrams \emph{cognitively privileged}
representations---recruiting spatial reasoning, free-ride inference, and
predictive processing---beyond being merely convenient notation? &
\autoref{sec:introduction}, \autoref{sec:cognitive-integration} &
\textbf{C6} \\
\textbf{RQ7} & Can the entire categorical framework be computationally
verified with real (non-mock) automated tests at ≥90\% coverage,
confirming that the abstractions are executable? &
\autoref{sec:diagrammatic-cognition}, \autoref{sec:daif-results} &
\textbf{C7} \\
\textbf{RQ8} & Do categorical string diagrams extend naturally into
quantum generalizations via TQNNs, ZX-calculus, and sheaf-theoretic
quantum semantics---and does quantum entanglement provide genuine
advantages for semantic communication? &
\autoref{sec:quantum-active-inference}, \autoref{sec:quantum-semantics}
& \textbf{C8} \\
\textbf{RQ9} & Can prompt injection attacks be formulated mathematically
as illicit cross-category morphisms (type violations), yielding a
\textbf{decidable} static check \textbf{relative to a fixed protocol /
interaction category} (compile-time where that grammar is enforced)? &
\autoref{sec:cognitive-security} & \textbf{C9} \\
\textbf{RQ10} & Does integrating enriched category theory with active
inference generate \emph{falsifiable} neurolinguistic predictions
(P600/N400 amplitudes scaling with enriched morphism weights)? &
\autoref{sec:cognitive-integration}, \autoref{sec:daif-results} &
\textbf{C10} \\
\textbf{RQ11} & Can grammatical case roles (NOM, ACC, INS) rigorously
type-check the tensor payloads of multi-agent AI framework protocols
(A2A, MCP, ANP) to prevent agentic hijacking? &
\autoref{sec:ai-implications} & \textbf{C11} \\
\bottomrule\noalign{}
\end{longtable}

\endgroup

\begin{quote}
\textbf{Reading guide.} \textbf{RQ1--RQ5} build the mathematical
framework cumulatively: each layer depends on the preceding one, from
case categories (RQ1) through string diagrams (RQ2), distributional
semantics (RQ3), enriched structure (RQ4), to topos-theoretic transfer
(RQ5). \textbf{RQ6} provides the cognitive science meta-argument that
threads through the entire paper, grounding diagrams in spatial
reasoning and predictive processing. \textbf{RQ7} supplies computational
validation of the framework via automated testing. \textbf{RQ8--RQ11}
extend the framework into four application domains: quantum semantics
(RQ8), cognitive security (RQ9), falsifiable neurolinguistic predictions
(RQ10), and multi-agent AI protocols (RQ11). The conclusion
(\autoref{sec:conclusion}) revisits each question with a summary of
results and identifies eight open directions (\textbf{F1}--\textbf{F8}).
\end{quote}
\end{frame}

\end{document}
